\section{Cyclic Groups and Cyclic Subgroups}

\begin{definition}[Cyclic]
    A group \( H \) is \textbf{cyclic} if \( H  \) can be generated by a single element; that is, there exists some element \( x \in H  \) such that \( H = \{ x^{n} : n \in \Z  \}  \) (where as usual the operation is multiplication).
\end{definition}

\begin{itemize}
    \item If the group in question is additive then we say that \( H \subseteq  G   \) is cyclic if 
        \[  H = \{ n x : n \in \Z  \}  \]
    \item We say that \( H  \) is generated by \( x  \) (or \( x  \) is the generator of \( H \)). 
    \item A cyclic group may have more than one generator.
\end{itemize}

\begin{prop}
   If \( H = \langle x \rangle   \), then \( | H  |  = | x  |  \) (where if one side of this equality is infinite, so is the other). More specifically, 
   \begin{enumerate}
       \item[(1)] If \( | H  |  = n < \infty  \), then \( x^{n} = 1  \) and \( 1, x, x^{2}, \dots, x^{n-1} \) are all distinct elements of \( H  \), and
        \item[(2)] If \( | H  |  = \infty  \), then \( x^{n} \neq 1  \) for all \( n \neq 0  \) and \( x^{a} \neq x^{b} \) for all \( a \neq b  \) in \( \Z  \).
   \end{enumerate}
\end{prop}

\begin{prop}[(3)]
   Let \( G  \) be an arbitrary group, \( x \in G  \) and let \( m,n \in \Z  \). If \( x^{n} = 1  \) and \( x^{m} = 1  \), then \( x^{d} = 1  \), where \( d = (m,n) \). In particular, if \( x^{m} = 1  \) for some \( m \in \Z  \), then \( | x  |  \mid m \). 
\end{prop}

\begin{theorem}[(4)]
   Any two cyclic groups of the same order are isomorphic. More specifically, 
   \begin{enumerate}
       \item[(1)] if \( n \in \Z^{+} \) and \( \langle x  \rangle\) and \( \langle y \rangle \) are both cyclic groups of order \( n  \), then the map 
           \begin{align*}
               \varphi : \langle x \rangle &\to \langle y \rangle  \\
                                         x^{k}  &\mapsto y^{k}
           \end{align*}
           is well defined and is an isomorphism.
   \end{enumerate}
\end{theorem}

\begin{prop}[(5)]
Let \( G  \) be a group, let \( x \in G  \) and let \( a \in \Z - \{  0  \}  \).
\begin{enumerate}
    \item[(1)] If \( | x  |  = \infty   \), then \( | x^{a} |  = \infty  \).
    \item[(2)] If \( | x  |  =  n < \infty  \), then \( | x^{a} |  = \displaystyle \frac{ n  }{ (n,a) }  \).
    \item[(3)] In particular, if \( | x  |  = n < \infty  \) and \( a  \) is a positive integer dividing \( n \), then \( | x^{a} |  = \frac{ n }{ a }  \).
\end{enumerate}
\end{prop}

\begin{prop}[(6)]
   Let \( H = \langle x  \rangle \). 
    \begin{enumerate}
        \item[(1)] Assume \( | x  |  = \infty  \). Then \( H = \langle x^{a} \rangle  \) if and only if \( a = \pm 1  \).
        \item[(2)] Assume \( | x  |  = n < \infty  \). Then \( H= \langle x^{a} \rangle\) if and only if \( (a,n) = 1   \). In particular, the number of generators of a \( H  \) is \( \varphi(n) \) (where \( \varphi  \) is Euler's \( \varphi \)-function).
    \end{enumerate}
\end{prop}

This is mainly used to give a description of all the subgroups of \( H  \) that generate \( H  \).

\begin{eg}[Multiplicative Group]
    Find all the subgroups of \( Z_{20} = \langle x  \rangle  \). 
    \begin{enumerate}
        \item[(1)] Find all the divisors of \( 20 \). That is, the divisors of \( 20  \) are 
            \[  \{ 1,2,4,5,10,20 \}. \]
        \item[(2)] From the above proposition (2) the divisors of \( 20  \) gives us the unique subgroups of \( \langle x  \rangle \) that generate \( Z_{20} \). These subgroups being 
            \[  \langle x \rangle, \langle x^{2} \rangle, \langle x^{4} \rangle, \langle x^{5} \rangle, \langle x^{10} \rangle, \langle x^{20} \rangle. \]
        \item[(3)] Using prop (5), we can find the order of each subgroup:
            \begin{align*}
                | \langle x^{2} \rangle | &= \frac{ 20 }{ (2,20) }  = \frac{ 20 }{ 2 }  = 10 \ \text{elements in} \ \langle x^{2} \rangle \\
                | \langle x^{4} \rangle |  &= \frac{ 20 }{ (4,20) } = \frac{ 20 }{ 4 } = 5 \ \text{elements in} \ \langle x^{4} \rangle \\
                | \langle x^{5} \rangle |  &= \frac{ 20 }{ (5,20) } = \frac{ 20 }{ 5 } = 4 \ \text{elements in} \ \langle x^{5} \rangle \\
                | \langle x^{10} \rangle |  &= \frac{ 20 }{ (10,20) } = \frac{ 20 }{ 10 } = 2 \ \text{elements in} \ \langle x^{10} \rangle \\
                | \langle x^{20} \rangle |  &= \frac{ 20 }{ (20,20) } = \frac{ 20 }{ 20 } = 1 \ \text{element in} \ \langle x^{20} \rangle \\
            \end{align*}
        \item[(4)] List out the generators of each subgroup by computing the Euler-Totient function \( \varphi(n) \). For example, given the cyclic subgroup \( \langle x^{2} \rangle \) which has order \( 10 \), we have 
            \[  \varphi(10) = \varphi(5 \cdot 2) = 4 \ \text{generators}.  \]
            Note that \( \varphi \) gives us the number of \( a \in \Z^{+} \) up to \( 10 \) such that \( (a,10) = 1  \) (where \( a  \) and \( 1 0 \) are relatively prime). Indeed, we have 
            \begin{itemize}
                \item \( (1,10) = 1 \implies (x^{2})^{1} = x^{2} \) is a generator;
                \item \( (2,10) = 2  \implies (x^{2})^{2} = x^{4} \) is NOT a generator;
                \item \( (3,10) = 1 \implies (x^{2})^{3} = x^{6} \) is a generator;
                \item \( (4,10) = 2 \implies (x^{2})^{4} = x^{8} \) is NOT a generator;
                \item \( (5,10) = 5 \implies (x^{2})^{5} = x^{10} \) is NOT a generator;
                \item \( (6,10) = 2 \implies (x^{2})^{6} = x^{12} \) is NOT a generator;
                \item \( (7,10) = 1 \implies (x^{2})^{7} = x^{14} \) is a generator;
                \item \( (8,10) = 2 \implies (x^{2})^{8} = x^{16} \) is NOT a generator;
                \item \( (9,10) = 1 \implies (x^{2})^{9} = x^{18} \) is a generator;
                \item \( (10,10) = 10 \implies (x^{2})^{10} = x^{20} \) is NOT a generator;
            \end{itemize}
        \item[(5)] To describe the containments of each of these subgroups in relation to each other, we fix each divisor in our list and identify which ones can be divided by it. For instance, we see that  
            \begin{align*}
                1 &\mid 2,4,5,10,20 \\
                2 &\mid 4,10,20 \\
                4 &\mid 20 \\
                5 &\mid 10,20 \\
                10 &\mid 20.
            \end{align*}
            This gives us the following containments:
            \begin{align*}
               \langle x^{20} \rangle &\subseteq \langle x^{10} \rangle \subseteq  \langle x^{5} \rangle \subseteq  \langle x  \rangle \\
               \langle x^{20} \rangle &\subseteq \langle x^{10} \rangle \subseteq \langle x  \rangle \\
               \langle x^{20} \rangle &\subseteq \langle x^{4} \rangle \subseteq  \langle x^{2}  \rangle \subseteq  \langle x \rangle \\
               \langle x^{20} \rangle &\subseteq \langle x^{2} \rangle \subseteq  \langle x \rangle \\
               \langle x^{20} \rangle &\subseteq \langle x  \rangle.
            \end{align*}
    \end{enumerate}
\end{eg}

\begin{theorem}[ ]
    Let \(H = \langle x  \rangle\) be a cyclic group.  
    \begin{enumerate}
        \item[(1)] Every subgroup of \( H  \) is cyclic. More precisely, if \( K \leq H  \), then either \( K = \{ 1  \}  \) or \( K = \langle x^{d} \rangle \), where \( d  \) is the smallest positive integer such that \( x^{d} \in K  \).
        \item[(2)] If \( | H  |  = \infty   \), then for any distinct nonnegative integers \( a \) and \( b  \), \( \langle x^{a} \rangle \neq \langle x^{b} \rangle  \). Furthermore, for every integer \( m  \), \( \langle x^{m} \rangle =  \langle x^{| m | } \rangle \), where \( | m |  \) denotes the absolute value of \( m  \), so that the nontrivial subgroups of \( H  \) correspond with the integers \( 1,2,3,\dots \). 
        \item[(3)] If \( | H  |  = n < \infty   \), then for each positive integer \( a  \) dividing \( n  \) there is a  unique subgroup of \( H  \) of order \( a  \). This subgroup is the cyclic group \( \langle x^{d} \rangle \), where \( d = \frac{ n }{ a }  \). Furthermore, for every integer \( m  \), \( \langle x^{m} \rangle= \langle x^{(n,m)} \rangle  \) so that the subgroups of \( H  \) correspond bijectively with the positive divisors of \( n  \).
    \end{enumerate}
\end{theorem}

